\subsection{Caratheodory-Hahn Extension}

\begin{definition} \label{def:premeasure-on-set-family} \hlt{Premeasures on Set Families}\\
Let $\emptyset\neq\mathcal{S}\subseteq\mathcal{P}(X)$ and $\mu:\mathcal{S}\to[0,+\infty]$.
\begin{enumerate}[label=(\roman*)]
\setlength{\itemsep}{0pt}
\item $\mu$ is \hlt{finitely additive on $\mathcal{S}$} if $\mu(\bigcup\limits_{k=1}^{n}E_k)=\sum\limits_{k=1}^{n}\mu(E_k)$ $\forall$ finite pairwise disjoint family $E_k\in\mathcal{S}$ with $\bigcup\limits_{k=1}^{n}E_k\in\mathcal{S}$.
\item $\mu$ is \hlt{countably monotone on $\mathcal{S}$} if, for all $E,E_1,E_2,\ldots\in\mathcal{S}$ with $E\subseteq\bigcup\limits_{k=1}^{\infty}E_k$, $\mu(E)\leq\sum\limits_{k=1}^{\infty}\mu(E_k)$.
\item $\mu$ is a \hlt{premeasure on $\mathcal{S}$} if $\mu$ is finitely additive on $\mathcal{S}$, countably monotone on $\mathcal{S}$, and $\emptyset\in\mathcal{S}\Rightarrow\mu(\emptyset)=0$.
\end{enumerate}
\hlt{Structural Commentary}: No closure property of $\mathcal{S}$ is assumed; additivity is imposed only for finite disjoint unions which remain in $\mathcal{S}$, and countable monotonicity only for countable $\mathcal{S}$-covers of members of $\mathcal{S}$.
\end{definition}

\begin{proposition} \label{prop:extension-forces-family-premeasure} \hlt{Necessary Conditions for Extension}\\
Let $\emptyset\neq\mathcal{S}\subseteq\mathcal{P}(X)$, let $\mu:\mathcal{S}\to[0,+\infty]$, let $\mu^\ast$ be the outer measure induced by $\mu$, and put $\overline{\mu}\defeq\mu^\ast|_{\mathcal{M}_{\mu^\ast}}$.\\
If $\mathcal{S}\subseteq\mathcal{M}_{\mu^\ast}$ and $\overline{\mu}(E)=\mu(E)$ for all $E\in\mathcal{S}$, then the following hold:
\begin{enumerate}[label=(\roman*)]
\setlength{\itemsep}{0pt}
\item $\emptyset\in\mathcal{S}\Rightarrow\mu(\emptyset)=0$.
\item $\mu$ is finitely additive on $\mathcal{S}$.
\item $\mu$ is countably monotone on $\mathcal{S}$.
\end{enumerate}
Hence $\mu$ is a premeasure on $\mathcal{S}$.
\begin{proof}
By the extension hypotheses, $\mu(F)=\overline{\mu}(F)=\mu^\ast(F)$ for every $F\in\mathcal{S}$.
\begin{enumerate}[label=(\roman*)]
\setlength{\itemsep}{0pt}
\item If $\emptyset\in\mathcal{S}$, then $\mu(\emptyset)=\overline{\mu}(\emptyset)=0$.
\item Let $n\geq1$, let $E_1,\ldots,E_n\in\mathcal{S}$ be pairwise disjoint, and assume $F\defeq\bigcup\limits_{k=1}^{n}E_k\in\mathcal{S}$.\\
Since $\overline{\mu}$ is a measure, $\mu(F)=\overline{\mu}(F)=\sum\limits_{k=1}^{n}\overline{\mu}(E_k)=\sum\limits_{k=1}^{n}\mu(E_k)$.
\item Let $E,E_1,E_2,\ldots\in\mathcal{S}$ and $E\subseteq\bigcup\limits_{k=1}^{\infty}E_k$. By monotonicity and countable subadditivity of $\mu^\ast$,\\
$\mu(E)=\mu^\ast(E)\leq\mu^\ast(\bigcup\limits_{k=1}^{\infty}E_k)\leq\sum\limits_{k=1}^{\infty}\mu^\ast(E_k)=\sum\limits_{k=1}^{\infty}\mu(E_k)$.
\end{enumerate}
\hlt{Structural Commentary}: The measure extension handles finite disjoint unions; $\mu^\ast$ handles countable covers.\\
\hlt{Key Insight}: Outer subadditivity supplies the inequality even when $\bigcup\limits_kE_k\notin\mathcal{S}$.
\end{proof}
\end{proposition}

\begin{definition} \label{def:relative-complements-and-semirings} \hlt{Relative Complements and Semirings}\\
Let $\emptyset\neq\mathcal{S}\subseteq\mathcal{P}(X)$.
\begin{enumerate}[label=(\roman*)]
\setlength{\itemsep}{0pt}
\item $\mathcal{S}$ is \hlt{closed under relative complements} if $A,B\in\mathcal{S}\Rightarrow A\setminus B\in\mathcal{S}$.
\item $\mathcal{S}$ is a \hlt{semiring of sets} if, for all $A,B\in\mathcal{S}$, $A\cap B\in\mathcal{S}$ and $A\setminus B=\bigsqcup\limits_{j=1}^{n}C_j$ for some $n\geq0$ and $C_j\in\mathcal{S}$.
\end{enumerate}
\hlt{Structural Commentary}: Relative-complement closure gives $\emptyset=A\setminus A\in\mathcal{S}$ and $A\cap B=A\setminus(A\setminus B)\in\mathcal{S}$.
\end{definition}

\begin{theorem} \label{thm:relative-complement-premeasure-extension} \hlt{Extension from Relative-Complement Families}\\
Let $\emptyset\neq\mathcal{S}\subseteq\mathcal{P}(X)$ be closed under relative complements, let $\mu:\mathcal{S}\to[0,+\infty]$ be a premeasure, and let $\mu^\ast$ be the outer measure induced by $\mu$. Then:
\begin{enumerate}[label=(\roman*)]
\setlength{\itemsep}{0pt}
\item $\mu^\ast(E)=\mu(E)$ for every $E\in\mathcal{S}$.
\item $\mathcal{S}\subseteq\mathcal{M}_{\mu^\ast}$.
\end{enumerate}
\begin{proof}
\begin{enumerate}[label=(\roman*)]
\setlength{\itemsep}{0pt}
\item Fix $E\in\mathcal{S}$. The one-set cover gives $\mu^\ast(E)\leq\mu(E)$. If $(E_k)_{k\geq1}\subseteq\mathcal{S}$ covers $E$, then countable monotonicity gives $\mu(E)\leq\sum\limits_{k=1}^{\infty}\mu(E_k)$; taking the infimum over all such covers gives $\mu(E)\leq\mu^\ast(E)$.
\item Fix $A\in\mathcal{S}$. By Lemma \ref{lem:caratheodory-finite-measure-test}, it suffices to prove $\mu^\ast(B\cap A)+\mu^\ast(B\setminus A)\leq\mu^\ast(B)$ for every $B\subseteq X$ with $\mu^\ast(B)<+\infty$. Let $\varepsilon>0$ and choose an $\mathcal{S}$-cover $(E_k)_{k\geq1}$ of $B$ with $\sum\limits_{k=1}^{\infty}\mu(E_k)<\mu^\ast(B)+\varepsilon$.\\
For each $k$, relative-complement closure gives $E_k\setminus A\in\mathcal{S}$ and $E_k\cap A=E_k\setminus(E_k\setminus A)\in\mathcal{S}$. By finite additivity, $\mu(E_k)=\mu(E_k\cap A)+\mu(E_k\setminus A)$. Since $(E_k\cap A)_{k\geq1}$ covers $B\cap A$ and $(E_k\setminus A)_{k\geq1}$ covers $B\setminus A$, $\mu^\ast(B\cap A)+\mu^\ast(B\setminus A)\leq\sum\limits_{k=1}^{\infty}\mu(E_k\cap A)+\sum\limits_{k=1}^{\infty}\mu(E_k\setminus A)=\sum\limits_{k=1}^{\infty}\mu(E_k)<\mu^\ast(B)+\varepsilon$. Let $\varepsilon\downarrow0$.
\end{enumerate}
\hlt{Structural Commentary}: Countable monotonicity identifies $\mu$ with $\mu^\ast$; finite additivity splits covers.\\
\hlt{Key Insight}: Relative complements turn one cover into inside/outside covers.
\end{proof}
\end{theorem}

\begin{remark} \label{rem:premeasure-roles-in-caratheodory-hahn} \hlt{Two Roles of a Premeasure}\\
In Theorem \ref{thm:relative-complement-premeasure-extension}, finite additivity is used to split each covering set across $A$ and $A^c$, while countable monotonicity is used to force $\mu^\ast(E)=\mu(E)$ on the original family.\\
\hlt{Key Insight}: Measurability comes from finite splitting; extension equality comes from countable cover control.
\end{remark}

\begin{proposition} \label{prop:semiring-premeasure-finite-union-extension} \hlt{Finite-Union Extension from a Semiring}\\
Let $\mathcal{S}$ be a semiring of subsets of $X$, and set $\mathcal{S}'\defeq\{\bigsqcup\limits_{k=1}^{n}A_k:n\geq0,\ A_1,\ldots,A_n\in\mathcal{S}\}$. If $\mu:\mathcal{S}\to[0,+\infty]$ is a premeasure, define $\mu':\mathcal{S}'\to[0,+\infty]$ by $\mu'(\bigsqcup\limits_{k=1}^{n}A_k)\defeq\sum\limits_{k=1}^{n}\mu(A_k)$ for any disjoint representation. Then:
\begin{enumerate}[label=(\roman*)]
\setlength{\itemsep}{0pt}
\item $\mathcal{S}'$ is closed under relative complements.
\item $\mu'$ is well-defined and extends $\mu$.
\item $\mu'$ is a premeasure on $\mathcal{S}'$.
\item $\mu'$ is the unique premeasure on $\mathcal{S}'$ extending $\mu$.
\end{enumerate}
\begin{proof}
\begin{enumerate}[label=(\roman*)]
\setlength{\itemsep}{0pt}
\item Let $E=\bigsqcup\limits_{i=1}^{r}A_i$ and $F=\bigsqcup\limits_{j=1}^{s}B_j$ in $\mathcal{S}'$.\\
For each $i$, repeated semiring differences decompose $A_i\setminus F$ into finitely many disjoint $\mathcal{S}$-sets.\\
These decompositions, over $i=1,\ldots,r$, form a finite disjoint $\mathcal{S}$-decomposition of $E\setminus F$, so $E\setminus F\in\mathcal{S}'$.
\item If $E=\bigsqcup\limits_{i=1}^{r}A_i=\bigsqcup\limits_{j=1}^{s}B_j$, then $A_i=\bigsqcup\limits_{j=1}^{s}(A_i\cap B_j)$ and $B_j=\bigsqcup\limits_{i=1}^{r}(A_i\cap B_j)$ with all intersections in $\mathcal{S}$.\\
Finite additivity of $\mu$ gives $\sum\limits_{i=1}^{r}\mu(A_i)=\sum\limits_{i=1}^{r}\sum\limits_{j=1}^{s}\mu(A_i\cap B_j)=\sum\limits_{j=1}^{s}\mu(B_j)$, so $\mu'$ is well-defined.\\
If $A\in\mathcal{S}$, the one-piece representation gives $\mu'(A)=\mu(A)$.
\item Since $\emptyset$ is the empty union, $\mu'(\emptyset)=0$.\\
Finite additivity of $\mu'$ follows by refining finitely many disjoint $\mathcal{S}'$-sets to their disjoint $\mathcal{S}$-pieces.\\
For countable monotonicity, let $E=\bigsqcup\limits_{i=1}^{r}A_i\in\mathcal{S}'$, let $F_m=\bigsqcup\limits_{j=1}^{N_m}B_{mj}\in\mathcal{S}'$, assume $E\subseteq\bigcup\limits_{m=1}^{\infty}F_m$.\\
For each $i$, $A_i\subseteq\bigcup\limits_{m,j}(A_i\cap B_{mj})$ with $A_i\cap B_{mj}\in\mathcal{S}$, so $\mu(A_i)\leq\sum\limits_{m,j}\mu(A_i\cap B_{mj})$.\\
Summing over $i$ gives $\mu'(E)\leq\sum\limits_m\sum\limits_j\sum\limits_i\mu(A_i\cap B_{mj})$.\\
For fixed $m,j$, repeated semiring differences decompose $B_{mj}\setminus\bigcup\limits_i(A_i\cap B_{mj})$ into finitely many disjoint $\mathcal{S}$-sets, so finite additivity gives $\sum\limits_i\mu(A_i\cap B_{mj})\leq\mu(B_{mj})$. Thus $\mu'(E)\leq\sum\limits_m\sum\limits_j\mu(B_{mj})=\sum\limits_m\mu'(F_m)$.
\item If $\nu$ is a premeasure on $\mathcal{S}'$ extending $\mu$, then finite additivity gives\\
$\nu(\bigsqcup\limits_{k=1}^{n}A_k)=\sum\limits_{k=1}^{n}\nu(A_k)=\sum\limits_{k=1}^{n}\mu(A_k)=\mu'(\bigsqcup\limits_{k=1}^{n}A_k)$ for every disjoint representation. Hence $\nu=\mu'$.
\end{enumerate}
\hlt{Structural Commentary}: Finite disjoint unions close a semiring under relative complements.\\
\hlt{Key Insight}: All estimates reduce to intersections of semiring pieces.
\end{proof}
\end{proposition}

\begin{theorem} \label{thm:caratheodory-hahn-semiring-extension} \hlt{Caratheodory-Hahn Extension from a Semiring}\\
Let $\mathcal{S}$ be a semiring of subsets of $X$, let $\mu:\mathcal{S}\to[0,+\infty]$ be a premeasure, let $\mu^\ast$ be the outer measure induced by $\mu$, and set $\overline{\mu}\defeq\mu^\ast|_{\mathcal{M}_{\mu^\ast}}$. Then:
\begin{enumerate}[label=(\roman*)]
\setlength{\itemsep}{0pt}
\item $\mathcal{S}\subseteq\mathcal{M}_{\mu^\ast}$ and $\overline{\mu}(A)=\mu(A)$ for every $A\in\mathcal{S}$.
\item If $X=\bigcup\limits_{k\geq1}S_k$ with $S_k\in\mathcal{S}$ and $\mu(S_k)<+\infty$, then $\overline{\mu}$ is $\sigma$-finite.
\item Under the hypothesis in (ii), if $\nu$ is a measure on $\mathcal{M}_{\mu^\ast}$ with $\nu(A)=\mu(A)$ for every $A\in\mathcal{S}$, then $\nu=\overline{\mu}$.
\end{enumerate}
\begin{proof}
\begin{enumerate}[label=(\roman*)]
\setlength{\itemsep}{0pt}
\item Let $(\mathcal{S}',\mu')$ be the finite-union extension from Proposition \ref{prop:semiring-premeasure-finite-union-extension}. The $\mathcal{S}$- and $\mathcal{S}'$-cover infima coincide, so $\mu^\ast$ is also induced by $\mu'$. Theorem \ref{thm:relative-complement-premeasure-extension} gives $\mathcal{S}'\subseteq\mathcal{M}_{\mu^\ast}$ and $\mu^\ast|_{\mathcal{S}'}=\mu'$. Since $\mu'|_{\mathcal{S}}=\mu$, the claim follows.
\item Assume $X=\bigcup\limits_{k\geq1}S_k$, $S_k\in\mathcal{S}$, $\mu(S_k)<+\infty$. By (i), each $S_k\in\mathcal{M}_{\mu^\ast}$, $\overline{\mu}(S_k)=\mu(S_k)<+\infty$, so $\overline{\mu}$ is $\sigma$-finite.
\item Assume (ii), and let $\nu$ be a measure on $\mathcal{M}_{\mu^\ast}$ extending $\mu$. By finite additivity, $\nu=\overline{\mu}=\mu'$ on $\mathcal{S}'$.\\
Hence every $\mathcal{S}'$-cover of $E\in\mathcal{M}_{\mu^\ast}$ gives $\nu(E)\leq$ its $\mu'$-mass, so $\nu(E)\leq\overline{\mu}(E)$. Put $H_n\defeq\bigcup\limits_{k=1}^{n}S_k\in\mathcal{S}'$.\\
Then $H_n\uparrow X$ and $\nu(H_n)=\overline{\mu}(H_n)=\mu'(H_n)<+\infty$.\\
Applying $\nu\leq\overline{\mu}$ to $E\cap H_n$ and $H_n\setminus E$ forces $\nu(E\cap H_n)=\overline{\mu}(E\cap H_n)$, because both measures split the same finite mass $H_n$. Since $E\cap H_n\uparrow E$, continuity from below gives $\nu(E)=\overline{\mu}(E)$.
\end{enumerate}
\hlt{Structural Commentary}: Pass to finite disjoint unions, then exhaust by finite-mass pieces.\\
\hlt{Key Insight}: $\mathcal{S}'$ supplies relative complements and the cover infimum for uniqueness.
\end{proof}
\end{theorem}

\begin{figure}[H]
  \centering
  \begin{tikzpicture}[
  every node/.style={font=\scriptsize},
  box/.style={draw=black!50, line width=0.6pt, align=center, text width=3.45cm, minimum height=0.70cm, inner sep=2pt, fill=white},
  source/.style={box, fill=blue!8},
  envelope/.style={box, fill=teal!8},
  measure/.style={box, fill=orange!12},
  arrow/.style={-{Latex[length=2mm]}, line width=0.55pt, black!65}
]
  \node[source] (mu) at (0,0) {$\mu:\mathcal{S}\to[0,+\infty]$\\premeasure on a semiring};
  \node[box] (prime) at (4.05,0) {$\mathcal{S}'$\\finite disjoint unions};
  \node[envelope] (mustar) at (8.10,0) {$\mu^\ast:2^X\to[0,+\infty]$\\induced outer measure};
  \node[measure] (mubar) at (12.15,0) {$\overline{\mu}=\mu^\ast|_{\mathcal{M}_{\mu^\ast}}$\\Caratheodory measure};

  \draw[arrow] (mu) -- (prime);
  \draw[arrow] (prime) -- (mustar);
  \draw[arrow] (mustar) -- (mubar);
\end{tikzpicture}

  \caption{Caratheodory-Hahn: pass to $\mathcal{S}'$, form $\mu^\ast$ by cover infima, then $\overline{\mu}\defeq\mu^\ast|_{\mathcal{M}_{\mu^\ast}}$ with $\overline{\mu}|_{\mathcal{S}}=\mu$.}
  \label{fig:caratheodory-hahn-extension-flow}
\end{figure}

\begin{corollary} \label{cor:sigma-finite-measures-agree-on-semiring} \hlt{Agreement on a $\sigma$-Finite Semiring}\\
Let $\mathcal{S}$ be a semiring of subsets of $X$, let $\mathcal{B}\defeq\sigma(\mathcal{S})$, and let $\nu_1,\nu_2$ be measures on $\mathcal{B}$ with $\nu_1|_{\mathcal{S}}=\nu_2|_{\mathcal{S}}$.\\
If $X=\bigcup\limits_{k\geq1}S_k$ for some $S_k\in\mathcal{S}$ with $\nu_1(S_k)=\nu_2(S_k)<+\infty$ for all $k$, then $\nu_1=\nu_2$.
\begin{proof} Let $\mathcal{S}'$ be the finite-disjoint-union closure of $\mathcal{S}$. Finite additivity gives $\nu_1=\nu_2$ on $\mathcal{S}'$.\\
Put $H_n\defeq\bigcup\limits_{k=1}^{n}S_k\in\mathcal{S}'$; then $H_n\uparrow X$ and $\nu_i(H_n)<+\infty$. For fixed $n$, $\mathcal{D}_n\defeq\{E\in\mathcal{B}:\nu_1(E\cap H_n)=\nu_2(E\cap H_n)\}$ is a Dynkin class containing the $\pi$-system $\mathcal{S}$, since $A\cap H_n\in\mathcal{S}'$ for $A\in\mathcal{S}$. Hence $\mathcal{D}_n=\mathcal{B}$.\\
Thus $\nu_1(E\cap H_n)=\nu_2(E\cap H_n)$ $\forall E\in\mathcal{B}$ and $n\geq1$; since $E\cap H_n\uparrow E$, continuity from below gives $\nu_1(E)=\nu_2(E)$.\\
\hlt{Structural Commentary}: Localize to finite-mass $H_n$, then apply $\pi$-$\lambda$.\\
\hlt{Key Insight}: Agreement on $\mathcal{S}$ propagates through $\mathcal{S}'$ and the exhaustion.
\end{proof}
\end{corollary}

\begin{example} \label{ex:bounded-interval-products-semirings} \hlt{Bounded Intervals and Rectangles Are Semirings}\\
Let $\mathcal{I}$ be the bounded intervals of $\R$, and let $\mathcal{R}_n\defeq\{I_1\times\cdots\times I_n:I_j\in\mathcal{I}\}$.
\begin{enumerate}[label=(\roman*)]
\setlength{\itemsep}{0pt}
\item $\mathcal{I}$ is a semiring: if $I,J\in\mathcal{I}$, then $I\cap J\in\mathcal{I}$ and $I\setminus J$ is a disjoint union of at most two members of $\mathcal{I}$.
\item $\mathcal{R}_n$ is a semiring in $\R^n$: intersections are coordinatewise products, and $R\setminus Q$ decomposes by coordinate cuts into finitely many disjoint rectangles in $\mathcal{R}_n$.
\end{enumerate}
Thus length on $\mathcal{I}$ and rectangular volume on $\mathcal{R}_n$ fit the Caratheodory-Hahn extension pattern.
\end{example}
